\documentclass[10pt,a4paper]{article}
\usepackage[T1]{fontenc}
\usepackage{lmodern,microtype}
\usepackage[margin=22mm]{geometry}
\usepackage{amsmath,graphicx,float}
\usepackage[numbers]{natbib}
\usepackage{xurl}
\usepackage{xcolor}
\usepackage[colorlinks=true,allcolors=blue!45!black]{hyperref}
\graphicspath{{../build/supplement/}{./}}
\setlength{\emergencystretch}{2em}
\renewcommand{\thefigure}{S2.\arabic{figure}}
\renewcommand{\thepage}{S2--\arabic{page}}
\begin{document}
\begin{center}
{\Large Supplement S2: Higher-dimensional M\"obius-strip layouts}\\[4pt]
Pawel Gajer and Jacques Ravel
\end{center}

These modern examples illustrate the higher-dimensional layout and projection
approach of \citet[][Section 3.6 and Figures 5--6]{gajer2004}. They use the
current \texttt{grip} backend, new graphs, and PCA, not the historical
implementation or projection.

\begin{figure}[H]
\centering
\includegraphics[width=\linewidth]{mobius-dimension-comparison.pdf}
\caption{Direct 3D layouts (top) and 4D layouts projected to 3D (bottom) of
unit-edge M\"obius-strip graphs with 150, 300, and 1,500 vertices. Each pair
uses the same graph, backend, settings, and seed. Thicker blue edges mark the
strip boundary. The lower row retains the first three principal components.
For display, configurations are centered, scaled to unit root-mean-square
radius, orthogonally aligned to a common reference, and viewed from one
direction. The projected twists appear somewhat smoother, but neither
appearance nor retained variance establishes improved graph-distance fidelity.}
\end{figure}

\paragraph{Construction and shared settings.}
A rectangular grid with $L$ longitudinal and $W$ transverse vertices is closed
by joining vertex $(L,j)$ to $(1,W+1-j)$, introducing a half-twist.
The grids $(L,W)=(30,5),(50,6),(100,15)$ have 270, 550, and 2,900 edges,
respectively, all assigned length one.
Both dimensions use \texttt{weighted.grip.nd()} with
\texttt{rounds = 160}, \texttt{final\_rounds = 256},
\texttt{num\_init = 24}, \texttt{num\_nbrs = 24}, and \texttt{seed = 1};
other arguments retain their defaults. Only this seed was used.
The analytic strip is used solely for display alignment (rotation or
reflection), never for layout initialization.

\paragraph{Projection and reproduction.}
PCA centers the 4D coordinates without scaling individual axes. Despite high
retained variance, projection can substantially shorten individual edges.
The supplied RDS records projected-to-original 4D edge-length ratios before
display normalization, coordinates, projection matrices, settings, and the R
session. The manuscript regenerates its 300-vertex pair using the same functions.
Reproduction commands are provided in \path{reproducibility/README.md}.

\bibliographystyle{plainnat}
{\small\bibliography{S2-mobius-comparison}}
\end{document}
